% ---------------------------------------------------------------
% ---------------------------------------------------------------
% Modelo de Trabalho Acadêmico utilizando classe repUERJ 
% desenvolvido no Instituto de Física da Universidade do Estado 
% do Rio de Janeiro (UERJ) para elaboração de teses, dissertação 
% e trabalhos de conclusão de curso.
%
% Este arquivo está editado na codificação de caracteres UTF-8.
%
% A classe repUERJ.cls foi criada a partir do código original 
% disponibilizado pelo grupo AbnTeX (coordenado por Gerald
% Weber). Foram feitas adequações para implementação das 
% normas de elaboração de teses e dissertações da UERJ.
%
% Os estilos repUERJformat.sty codificam os elementos
% pré-textuais e pós-textuais pelas nornas de elaboração de
% dissertações e teses da UERJ.
%
% O estilo repUERJpseudocode.sty codifica a elaboração de
% algoritmos utilizando um glossário diferente dos pacotes 
% tradicionais (usa pseudocomando em portugues).
%
% As normas da UERJ para elaboração de teses e dissertações 
% pode ser obtidas no documento disponível no site
%      http://www.bdtd.uerj.br/roteiro_uerj_web.pdf
%
% Agradecimentos ao NPROTEC/Rede Sirius/UERJ e à Biblioteca
% Setorial da Física.
%
% As referencia estão baseadas em bibtex e citação no sistema 
% autor-data, podendo ser alterada para o sistema numerado.
%
% GUIA DE USO DO PACOTE repUERJ SE ENCONTRA NO DIRETÓRIO DOCUMENTOS.
%
% ---------------------------------------------------------------
% ---------------------------------------------------------------
%
% MODELO ATUALIZADO EM 2023/12/13
%
% ---------------------------------------------------------------
% ---------------------------------------------------------------
%
\documentclass[a4paper,12pt,oneside,onecolumn,final,fleqn]{repUERJ}
% ---
% OBRIGATÓRIO
% Pacotes fundamentais 
% ---
\usepackage[brazil]{babel}  % adequação para o português Brasil
\usepackage[utf8]{inputenc} % Determina a codificação utilizada
                            % (conversão automática dos acentos)
\usepackage{indentfirst}    % Indenta o primeiro paragrafo de
                            % cada seção.
\usepackage{graphicx}       % Inclusão de gráficos
\usepackage{subfig}
\usepackage{multirow}
\usepackage{amsmath,amssymb}  % pacotes matemáticos
\usepackage[maxfloats=25]{morefloats}
\usepackage{array}

\usepackage{makeidx}        % Cria o índice
\usepackage{hyperref}       % Controla a formação do índice

% ----------------------------------------------------------------
% OBRIGATÓRIO
% estilo que implementa os comandos principais do pacote repUERJ
% pacote para as normas da UERJ
% ----------------------------------------------------------------
\usepackage[font=default,frame=no]{repUERJformat} 

% ----------------------------------------------------------------
% para referências no estilo autor-data
% descomente a linha 72 e comente as linhas 78-80
% ----------------------------------------------------------------
\usepackage[alf]{abntex2cite} % pacote para citacoes autor-data

% ----------------------------------------------------------------
% para referências no estilo numerado
% descomente as linhas 78-80 e comente a linha 72
% ----------------------------------------------------------------
%\usepackage[num]{abntex2cite} % pacote para citacoes com numeracao
%\citebrackets[]
%\setlength{\bibindent}{1.5em}

% ----
\citeoption{abnt-options4}
% ----

% ----------------------------------------------------------------
% Retira os escapes (chavrons) < e > das urls das referências 
% bibliográficas
% ----------------------------------------------------------------
\def\UrlLeft{}\def\UrlRight{}
% ----------------------------------------------------------------

% ----
\setlength\extrarowheight{2pt}
% ----

% ----------------------------------------------------------------
% OPCIONAL
% pacote auxiliar para elaboração de pseudocódigos
% este pacote pode ser retirado caso nao se planeje
% elaborar pseudocódigos
% ----------------------------------------------------------------
\usepackage[dots=yes,num=yes]{repUERJpseudocode}

% % ----------------------------------------------------------------
% OPCIONAL
% % Este trecho de comandos pode ser apagado
% %    para gerar a versão final sem as referências retroativas
% % ----------------------------------------------------------------
% \usepackage[hyperpageref,brazil]{backref}
% % Usado sem a opção hyperpageref de backref
% \renewcommand{\backrefpagesname}{>>> Citado na(s) página(s):~}
% % Texto padrão antes do número das páginas
% \renewcommand*{\backref}{}
% % Define os textos da citação
% \renewcommand*{\backrefalt}[4]{
% \ifcase #1 %
% [> Nenhuma citação no texto. <]%
% \or
% [> Citado na página #2. <]%
% \else
% [> Citado #1 vezes nas páginas #2. <]%
% \fi}%

% ----------------------------------------------------------------
% OPCIONAL
% verbatim para os campos do bibtex no apêndice
% configura a apresentação dos exemplos no apêndice
% pode ser retirado!!!
% ----------------------------------------------------------------
\usepackage{framed,color,verbatim}
\definecolor{shadecolor}{rgb}{.95, .95, .95}
\newenvironment{vbrcode}%
   {\snugshade\verbatim}%
   {\endverbatim\endsnugshade}



% ***************************************************************
% Informações de autoria e institucionais
% ***************************************************************
% Ler o guia de uso deste pacote para melhor aproveitamento. 
% Ver o documento “Guia de uso (repUERJ).pdf” no diretório “documentos”.
% ***************************************************************

%----------------------------------------------------------------
% Imagens pretextuais (precisam estar no mesmo diretório deste arquivo .tex)
%----------------------------------------------------------------

\logo{imagens/logo_uerj_cinza.png}
\marcadagua{imagens/marcadagua_uerj_cinza.png}{1}{160}{255}

%----------------------------------------------------------------
% Informações da instituição
%----------------------------------------------------------------

\instituicao{Universidade do Estado do Rio de Janeiro}
            {Centro setorial}  
            {Unidade}{Patrono} 

%----------------------------------------------------------------
% Informações da autoria do documento
%----------------------------------------------------------------

% \oautor{Nóme}{Sóbrenome}
%       {Iniciais do nome} % iniciais do nome

\autor{Nome}{Sobrenome}
      {Iniciais do nome} % iniciais do nome

\titulo{Título do trabalho}
\title{Título do trabalho em inglês}

% se não for usar a quarta palavra chave, deixar o campo vazio: {}
\palavraschaves{Primeira palavra-chave}
               {Segunda palavra-chave}
               {Terceira palavra-chave}
               {Quarta palavra-chave (opcional) ou vazio}

\keywords{First keyword}
         {Second keyword}
         {Third keyword}
         {Fourth keyword or empty}

\orientador{Cargo Titulação} 
           {Nome}{Sobrenome} 
           {Unidade -- UERJ} 

\coorientador{Cargo Titulação} 
             {Nome}{Sobrenome} 
             {Unidade -- Instituição} 

%----------------------------------------------------------------
% Grau pretendido (Doutor, Mestre, Bacharel, Licenciado) e Curso
%----------------------------------------------------------------

\grau{Graduado} % Doutor, Mestre, Bacharel, Licenciado
\curso{Curso}
%\areadeconcentracao{linha de pesquisa} % opcional

%----------------------------------------------------------------
% Informações adicionais (local, data e paginas)
%----------------------------------------------------------------

\local{Rio de Janeiro} 
\data{dd}{mês}{aaaa} 

% ***************************************************************
% Configurações de aparência do PDF final
% ***************************************************************

% alterando o aspecto da cor azul
%\definecolor{blue}{RGB}{41,5,195}
%\definecolor{apricot}{RGB}{251,206,177}

% informações do PDF
\hypersetup{
  unicode=false,
  pdftitle={\UERJtitulo},
  pdfauthor={\UERJautor},
  pdfsubject={\UERJpreambulo},
  pdfkeywords={PALAVRAS-CHAVES:}{ \chaveA}{ \chaveB}{ \chaveC}{ \chaveD},
  pdfproducer={\packagename}, % producer of the document
  % para tirar as cores, ponha tudo em preto (black)
  colorlinks=true,    % false: boxed links; true: colored links
  linkcolor=blue,     % color of internal links blue
  citecolor=purple,   % color of links to bibliography blue
  filecolor=red,      % color of file links magenta
  urlcolor=magenta,
  bookmarksdepth=4,
%   backref=true,
%   pagebackref=true,
%   bookmarks=true,
}

% ***************************************************************
% Índice remissivo
% ***************************************************************
%
%\makeindex % compila o índice; se não for usar, comentar
%
% ***************************************************************
% Fim do preâmbulo
% ***************************************************************

%/\/\/\/\/\/\/\/\/\/\/\/\/\/\/\/\/\/\/\/\/\/\/\/\/\/\/\/\/\/
\begin{document}
%/\/\/\/\/\/\/\/\/\/\/\/\/\/\/\/\/\/\/\/\/\/\/\/\/\/\/\/\/\/

%XXXXXXXXXXXXXXXXXXXXXXXXXXXXXXXXXXXXXXXXXXXXXXXXXXXXXXXXXXXXXXXX
% ELEMENTOS PRE-TEXTUAIS
%XXXXXXXXXXXXXXXXXXXXXXXXXXXXXXXXXXXXXXXXXXXXXXXXXXXXXXXXXXXXXXXX
\frontmatter % inicia a área dos elementos pré-textuais
%XXXXXXXXXXXXXXXXXXXXXXXXXXXXXXXXXXXXXXXXXXXXXXXXXXXXXXXXXXXXXXXX

% ----------------------------------------------------------
% Capa e a folha de rosto
% ----------------------------------------------------------
%
\capa
\folhaderosto
%
% ----------------------------------------------------------
% Inserir a ficha catalográfica
% ----------------------------------------------------------
% ---
% A biblioteca deverá providenciar a ficha catalográfica. 
% Salve a ficha no formato PDF. Use o nome do arquivo PDF 
% como argumento do comando. 
%
% Exemplo: ficha catalográfica no arquivo 'ficha.pdf'
%     \fichacatalografica{ficha.pdf}
%
% Enquanto não possuir a ficha catalográfica, use o comando sem
% argumentos.
% ---
%
\fichacatalografica{}
%
% ----------------------------------------------------------
% Folha de aprovação
% ----------------------------------------------------------
%
\begin{folhadeaprovacao}
  \assinatura{Cargo Título Nome Completo}
             {Instituição por extenso}
  \assinatura{Cargo Título Nome Completo}
             {Instituição por extenso}
  \assinatura{Cargo Título Nome Completo}
             {Instituição por extenso}
  \assinatura{Cargo Título Nome Completo}
             {\UERJunidade \UERJunidadenome\ -- UERJ}
  \assinatura{Cargo Título Nome Completo}
             {\UERJunidade \UERJunidadenome\ -- UERJ}
  \assinatura{Cargo Título Nome Completo}
             {\UERJunidade \UERJunidadenome\ -- UERJ}
\end{folhadeaprovacao}
%
% ----------------------------------------------------------
% Dedicatória
% ----------------------------------------------------------
%
\pretextualchapter{Dedicatória}
\vfill
Texto da dedicatória (opcional).
%
% ----------------------------------------------------------
% Agradecimentos
% ----------------------------------------------------------
%
\pretextualchapter{Agradecimentos}

Texto de agradecimento (opcional).
%
% ----------------------------------------------------------
% Epigrafe (opcional)
% ----------------------------------------------------------
%
\pretextualchapter{}
  \vfill
  \begin{flushright}
 (opcional)\\
 Texto da epígrafe\\
 Pensamento, reflexão\\    
    \textit{autor}
  \end{flushright}
%
% ----------------------------------------------------------
% RESUMO
% ----------------------------------------------------------
%
\pretextualchapter{Resumo}
\referencia % linha em branco depois

Texto do resumo em português.
~\\

\imprimirchaves % linha em branco antes
%
% ----------------------------------------------------------
% Abstract
% ----------------------------------------------------------
%
\pretextualchapter{Abstract}
\reference % linha em branco depois

Texto do resumo em inglês.
~\\

\printkeys % linha em branco antes
%
% ----------------------------------------------------------
% Listas de ilustrações e tabelas
% ----------------------------------------------------------
%
\listadefiguras    %\begin{figure}{largura}...\end{figure}
\listadegraficos   %\begin{graph}{largura}...\end{graph}
\listadetabelas    %\begin{table}{largura}...\end{table}
\listadequadros    %\begin{inframe}{largura}...\end{inframe}
%
% ----------------------------------------------------------
% Outras listas
% ----------------------------------------------------------
%
\listadealgoritmos % opcional %\begin{algorithm}...\end{algorithm}
%
% ----------------------------------------------------------
% Lista de abreviaturas e siglas (opcional)
% ----------------------------------------------------------
%
\pretextualchapter{Lista de abreviaturas e siglas}
    \abreviatura{sigla 1}{por extenso}
    \abreviatura{sigla 2}{por extenso}
    \abreviatura{sigla 3}{por extenso}
%
% ----------------------------------------------------------
% Lista de simbolos (opcional)
% ----------------------------------------------------------
%
\pretextualchapter{Lista de símbolos}
    \simbolo{$simbolo 1$}{significado e/ou valor}
    \simbolo{$simbolo 2$}{significado e/ou valor}
    \simbolo{$simbolo 3$}{significado e/ou valor}
%
% ----------------------------------------------------------
% Sumario
% ----------------------------------------------------------
%
\sumario
%
%XXXXXXXXXXXXXXXXXXXXXXXXXXXXXXXXXXXXXXXXXXXXXXXXXXXXXXXXXXXXXXXX
% ELEMENTOS TEXTUAIS
%XXXXXXXXXXXXXXXXXXXXXXXXXXXXXXXXXXXXXXXXXXXXXXXXXXXXXXXXXXXXXXXX
\mainmatter % inicia a área de desenvolvimento do conteúdo
%XXXXXXXXXXXXXXXXXXXXXXXXXXXXXXXXXXXXXXXXXXXXXXXXXXXXXXXXXXXXXXXX

%================================================================
\chapter*{Introdução}
%================================================================

\begin{epigrafeonline}
\hfill Texto da epígrafe \textit{in locu} (opcional).\\
\hspace*{\fill}\textit{Autor}\\
\end{epigrafeonline}

Texto da introdução.

\nocite{bib:ABNT2011}

\nocite{bib:livro}

\nocite{bib:ebook}

\nocite{bib:tese}

\nocite{bib:dissertacao}

\nocite{bib:monografia}

\nocite{bib:TCC}

\nocite{bib:artigo}

\nocite{bib:proceedings}

\nocite{bib:inproceedings}

\nocite{bib:relatorio}

\nocite{bib:software}

\nocite{bib:software1}

\nocite{bib:inbook}

\nocite{bib:incollection}


%================================================================
\chapter{T\'itulo do cap\'itulo}
%================================================================

\begin{epigrafeonline}
\hfill Texto da epígrafe \textit{in locu} (opcional).\\
\hspace*{\fill}\textit{Autor}\\
\end{epigrafeonline}

Texto do capítulo.

%----------------------------------------------------------------
\section{Título da seção}
%----------------------------------------------------------------

Texto da seção.

%----------------------------------------------------------------
\subsection{Título da subseção}
%----------------------------------------------------------------

Texto da subseção

%----------------------------------------------------------------
\subsubsection{Título da subsubseção}
%----------------------------------------------------------------

Texto da subsubseção

%================================================================
\chapter*{Conclusão}
%================================================================

Texto da conclusão.

%\index{Introdução!Capítulo}.

%XXXXXXXXXXXXXXXXXXXXXXXXXXXXXXXXXXXXXXXXXXXXXXXXXXXXXXXXXXXXXXXX
% ELEMENTOS POS-TEXTUAIS
%XXXXXXXXXXXXXXXXXXXXXXXXXXXXXXXXXXXXXXXXXXXXXXXXXXXXXXXXXXXXXXXX
\backmatter % inicia a área dos elementos pós-textuais
%XXXXXXXXXXXXXXXXXXXXXXXXXXXXXXXXXXXXXXXXXXXXXXXXXXXXXXXXXXXXXXXX

%===========================================================
% Referencias via BibTeX
%===========================================================
\bibliography{bibliografia,exemplos}

%===========================================================
\postextualchapter*{Glossário} % elemento opcional
%===========================================================
\definicao{termo 1}{significado}
\definicao{termo 2}{significado}
\definicao{termo 3}{significado}

%XXXXXXXXXXXXXXXXXXXXXXXXXXXXXXXXXXXXXXXXXXXXXXXXXXXXXXXXXXX
% OPCIONAL
% Apêndice
%XXXXXXXXXXXXXXXXXXXXXXXXXXXXXXXXXXXXXXXXXXXXXXXXXXXXXXXXXXX
\appendix % inicia os apêndices
%XXXXXXXXXXXXXXXXXXXXXXXXXXXXXXXXXXXXXXXXXXXXXXXXXXXXXXXXXXX

%===========================================================
\postextualchapter{Título do apêndice}
%===========================================================

% ----------------------------------------------------------
\section{Título da seção}
% ----------------------------------------------------------

Texto da seção.

% ----------------------------------------------------------
\subsection{Título da subseção}
% ----------------------------------------------------------

Texto da subseção.

% ----------------------------------------------------------
\subsubsection{Título da subsubseção}
% ----------------------------------------------------------

Texto da subsubseção.

%===========================================================
\postextualchapter{Título do apêndice}
%===========================================================

% ----------------------------------------------------------
\section{Título da seção}
% ----------------------------------------------------------

Texto da seção.

% ----------------------------------------------------------
\subsection{Título da subseção}
% ----------------------------------------------------------

Texto da subseção.

% ----------------------------------------------------------
\subsubsection{Título da subsubseção}
% ----------------------------------------------------------

Texto da subsubseção.

%XXXXXXXXXXXXXXXXXXXXXXXXXXXXXXXXXXXXXXXXXXXXXXXXXXXXXXXXXXX
% OPCIONAL
% Anexo
%XXXXXXXXXXXXXXXXXXXXXXXXXXXXXXXXXXXXXXXXXXXXXXXXXXXXXXXXXXX
\annex % inicia os anexos
%XXXXXXXXXXXXXXXXXXXXXXXXXXXXXXXXXXXXXXXXXXXXXXXXXXXXXXXXXXX

%===========================================================
\postextualchapter{Título do anexo}
%===========================================================

% ----------------------------------------------------------
\section{Título da seção}
% ----------------------------------------------------------

Texto da seção.

% ----------------------------------------------------------
\subsection{Título da subseção}
% ----------------------------------------------------------

Texto da subseção.

% ----------------------------------------------------------
\subsubsection{Título da subsubseção}
% ----------------------------------------------------------

Texto da subsubseção.

%===========================================================
\postextualchapter{Título do anexo}
%===========================================================

% ----------------------------------------------------------
\section{Título da seção}
% ----------------------------------------------------------

Texto da seção.

% ----------------------------------------------------------
\subsection{Título da subseção}
% ----------------------------------------------------------

Texto da subseção.

% ----------------------------------------------------------
\subsubsection{Título da subsubseção}
% ----------------------------------------------------------

Texto da subsubseção.

%/\/\/\/\/\/\/\/\/\/\/\/\/\/\/\/\/\/\/\/\/\/\/\/\/\/\/\/\/\/
\end{document}
%/\/\/\/\/\/\/\/\/\/\/\/\/\/\/\/\/\/\/\/\/\/\/\/\/\/\/\/\/\/
